\newpage

% ######################################################################
% PART IV: DEEP DIVES
% ######################################################################
\part{Deep Dives --- Andar Ki Duniya}

% ######################################################################
% SECTION 21: LANGGRAPH INTERNALS
% ######################################################################
\section{LangGraph Internals --- Graph Engine Ki Andar Ki Baat}

\subsection{StateGraph Kya Hai, Asal Mein?}

\begin{storybox}
Ab tak humne \texttt{graph/workflow.py} ko line-by-line dekha. Ab is section mein hum \textbf{andar jayenge} --- LangGraph ka engine kaise kaam karta hai, node kab chalta hai, state kaise flow hoti hai, aur conditional edge ka jadoo kya hai. Ye concepts interview mein \textbf{gold} hain --- kyunki har koi bolta hai ``maine LangGraph use kiya'', par koi nahi batata ki uske andar kya ho raha hai.
\end{storybox}

\begin{conceptbox}
\textbf{Graph Abstraction ka Core Idea:}

Ek LangGraph app = \textbf{state} + \textbf{nodes} + \textbf{edges}.

\begin{itemize}
\item \textbf{State:} Ek TypedDict (\texttt{ResearchState}) jo pura pipeline ki memory hai --- topic, research\_data, analysis, report, sab kuch.
\item \textbf{Nodes:} Functions jo state lete hain aur ek partial update return karte hain.
\item \textbf{Edges:} Bataate hain ki kaun sa node kis ke baad chalega.
\item \textbf{Conditional Edges:} Runtime par decide karte hain ki agla node kaun sa hai.
\end{itemize}

Iska MATLAB: pipeline ka flow code mein hardcoded nahi hai --- wo \textbf{graph definition} hai jo runtime par execute hota hai.
\end{conceptbox}

\subsection{Node Execution Ka Mechanism}

\begin{codestorybox}
\textbf{Node function ki signature (hamesha):}

\begin{lstlisting}[style=pythonstyle, caption={Ek node ka contract}]
def researcher_node(state: ResearchState) -> dict:
    # state = ab tak ka pura data (read-only view)
    # kuch kaam karo
    return {
        "research_data": all_results,
        "messages": [...],
    }
\end{lstlisting}

\textbf{Contract:}
\begin{itemize}
\item Input: poora \texttt{state} dictionary (read karo, mutate mat karo).
\item Output: sirf \textbf{jo fields update karne hain} unka dict.
\item LangGraph output dict ko purani state par \textbf{merge} karta hai.
\end{itemize}
\end{codestorybox}

\begin{mathbox}
\textbf{State update ka math:}

Agar purani state $S$ hai aur node $\delta$ update return karta hai, to:
\[
S' = S \oplus \delta
\]
Jahan $\oplus$ default merge = shallow dictionary merge. Special fields (jaise \texttt{messages} with \texttt{add\_messages} reducer) ka apna merge rule hota hai.

\textbf{Shallow merge MATLAB:} nested values replace ho jaati hain, top-level keys add/update hoti hain. Isliye har node sirf apne fields return karta hai --- baaki state safe rehti hai.
\end{mathbox}

\subsection{add\_messages Reducer --- Messages Ka Jadoo}

\begin{conceptbox}
\texttt{messages} field par \texttt{Annotated[list, add\_messages]} lagaya hai. Ye LangGraph ko batata hai: \textbf{``is field ko simple replace mat karo, balki append/merge karo.''}

Agar har node \texttt{messages: ["naya message"]} return kare, to:
\begin{itemize}
\item Normal merge: pichla message \textbf{erase} ho jata (sirf last wala bachega).
\item \texttt{add\_messages} merge: naya message \textbf{append} hota hai.
\end{itemize}

Is project mein actually \texttt{messages} ko sirf logging ke liye use kiya hai --- par architecture mein ye pattern important hai kyunki RAG/chat agents mein message history accumulate karni hoti hai.
\end{conceptbox}

\subsection{Compile() Ke Andar Kya Hota Hai}

\begin{codestorybox}
\begin{lstlisting}[style=pythonstyle]
graph = StateGraph(ResearchState)
graph.add_node("researcher", researcher_node)
graph.add_node("analyst", analyst_node)
graph.add_node("fact_checker", fact_checker_node)
graph.add_node("writer", writer_node)
graph.add_edge(START, "researcher")
graph.add_edge("researcher", "analyst")
graph.add_edge("analyst", "fact_checker")
graph.add_conditional_edges(
    "fact_checker",
    should_continue,
    {"continue": "researcher", "end": "writer"}
)
graph.add_edge("writer", END)
research_graph = graph.compile()
\end{lstlisting}

\textbf{compile() kya karta hai:}
\begin{itemize}
\item Nodes aur edges ko ek internal \textbf{execution plan} mein convert karta hai.
\item Har node ke aage/peeche validation layers add karta hai.
\item Ek \textbf{CompiledStateGraph} object deta hai jisme \texttt{.invoke()}, \texttt{.stream()}, \texttt{.get\_graph()} jaise methods hain.
\item Compile par hi LangGraph ko pata chalta hai ki graph \textbf{connected} hai ya nahi (har node reachable hai?).
\end{itemize}
\end{codestorybox}

\subsection{invoke() vs stream() --- Do Tarike}

\begin{notebox}
\textbf{invoke()}: Poora graph end-to-end chalao, final state lao. Blocking call.

\textbf{stream()}: Har node ke baad intermediate results yield karo. Real-time progress dikhane ke liye (Streamlit mein use hota hai).

\begin{itemize}
\item \texttt{invoke} returns \texttt{dict} --- final state.
\item \texttt{stream} returns generator of \texttt{(node\_name, output)} tuples.
\end{itemize}
\end{notebox}

\begin{codestorybox}
\textbf{Streamlit app mein dono ka use:}

\begin{lstlisting}[style=pythonstyle, caption={stream vs invoke in practice}]
# Streamlit: real-time streaming
for event in research_graph.stream(state, config=config):
    for node_name, node_output in event.items():
        st.session_state.agent_statuses[node_name] = "done"

# Django / MCP: final result chahiye
final_state = research_graph.invoke(state)
report = final_state["final_report"]
\end{lstlisting}
\end{codestorybox}

\subsection{should\_continue --- Graph Ka Traffic Cop}

\begin{codestorybox}
\begin{lstlisting}[style=pythonstyle, caption={Conditional edge logic}]
def should_continue(state: ResearchState) -> str:
    if state.get("fact_check_passed", False):
        return "end"
    if state.get("revision_count", 0) >= MAX_REVISIONS:
        return "end"
    return "continue"
\end{lstlisting}

\textbf{Traffic cop ka kaam:}
\begin{itemize}
\item \texttt{fact\_check\_passed == True} $\rightarrow$ \texttt{"end"} $\rightarrow$ writer chalao.
\item Revision count limit cross ho gaya $\rightarrow$ \texttt{"end"} (infinite loop se bachao).
\item Warna \texttt{"continue"} $\rightarrow$ researcher se naye queries ke saath dubara research.
\end{itemize}

\textbf{Why this matters in interviews:} Ye ek \textbf{feedback loop} hai --- agent apne output ko khud verify karta hai aur improve karta hai. Single-prompt LLM mein ye loop nahi hota.
\end{codestorybox}

\begin{warningbox}
\textbf{Infinite loop ka danger:} Agar \texttt{MAX\_REVISIONS} check nahi hota aur fact\_checker hamesha FAIL deta, to graph \texttt{researcher $\rightarrow$ analyst $\rightarrow$ fact\_checker} mein kabhi na khatam hone wala chakkar laga deta. Isliye:
\begin{enumerate}[label=\arabic*.]
\item \texttt{revision\_count} har loop mein increment hota hai.
\item \texttt{MAX\_REVISIONS} (default 3) par hard stop.
\item Stop hone par bhi report generate hoti hai (best-effort).
\end{enumerate}
\end{warningbox}

\subsection{LangGraph Bina --- Manual Loop Kaisa Dikhta}

\begin{codestorybox}
\textbf{Agar LangGraph na hota, to pipeline aise likhni padti:}

\begin{lstlisting}[style=pythonstyle, caption={Manual pipeline (no framework)}]
def run_pipeline(topic):
    state = initial_state(topic)
    for revision in range(MAX_REVISIONS + 1):
        state.update(researcher_node(state))
        state.update(analyst_node(state))
        result = fact_checker_node(state)
        state.update(result)
        if result["fact_check_passed"]:
            break
    state.update(writer_node(state))
    return state
\end{lstlisting}

\textbf{Problems manual loop mein:}
\begin{itemize}
\item Flow logic (kaun kis ke baad) code mein khoya hai --- ek \texttt{for} loop ke andar.
\item Naya node add karna = manual loop ko todna.
\item Streaming / checkpointing / parallel nodes khud banana padega.
\item Testing hard --- pura pipeline ek monolithic function.
\end{itemize}

\textbf{LangGraph gives:} declarative definition, built-in streaming, checkpointing, visualization, parallelism, aur clean testing.
\end{codestorybox}

\begin{interviewbox}
\textbf{Q: ``StateGraph vs simple function chain --- overhead kya hai?''}

\textbf{Answer:} Honest answer --- is project ke scale par (4 nodes, linear flow) framework ka overhead minimal hai, aur benefits practical hain: (1) Conditional loop cleanly express hota hai; (2) stream() se UI progress milta hai bina alag threading ke; (3) Agar future mein parallel nodes (jaise 2 researchers alag topics par) chahiye, graph extend karna trivial hai. Overhead sirf ek extra abstraction layer hai jo code readability ko improve karta hai.
\end{interviewbox}

\subsection{LangGraph vs CrewAI vs AutoGen --- Comparison}

\begin{longtable}{p{3.2cm} p{3.4cm} p{3.4cm} p{3.4cm}}
\toprule
\textbf{Feature} & \textbf{LangGraph} & \textbf{CrewAI} & \textbf{AutoGen} \\
\midrule
\textbf{Core model} & Graph/state machine & Role-based crews & Conversational agents \\
\textbf{Control flow} & Explicit edges, conditions & Sequential/parallel tasks & Agent-to-agent chat \\
\textbf{State management} & TypedDict + reducers & Task outputs & Conversation history \\
\textbf{Streaming} & Built-in (.stream) & Limited & Built-in \\
\textbf{Debugging} & Graph visualization, tracing & Medium & Medium \\
\textbf{Learning curve} & Steep (graph thinking) & Gentle & Medium \\
\textbf{Is project mein?} & \textbf{YES} --- full control & No & No \\
\bottomrule
\end{longtable}

\begin{conceptbox}
\textbf{Is project mein LangGraph kyun?} Kyunki is pipeline ka \textbf{heart} ek \textbf{conditional feedback loop} hai (fact-check $\rightarrow$ revise). Graph model isko naturally express karta hai. CrewAI role-based approach mein aisa loop force karna awkward hota. AutoGen conversational hai --- research pipelines ke liye less structured.
\end{conceptbox}

% ######################################################################
% SECTION 22: EMBEDDINGS DEEP DIVE
% ######################################################################
\section{Embeddings Deep Dive --- Vector Banane Ka Science}

\subsection{Embedding Kya Hai --- Ek Sentence Se 384 Numbers Tak}

\begin{storybox}
Humne \texttt{rag/vector\_store.py} mein \texttt{HuggingFaceEmbeddings} use kiya. Ab samjho: jab hum PDF ka ek chunk lete hain, \texttt{all-MiniLM-L6-v2} model us chunk ko \textbf{384 numbers} ki ek list mein convert karta hai. Ye numbers = us text ka \textbf{semantic fingerprint}.

\textbf{Key insight:} Do similar sentences ke embeddings \textbf{near} hote hain (vector space mein distance kam), dissimilar ke far. Isi par similarity search chalti hai.
\end{storybox}

\begin{mathbox}
\textbf{Cosine Similarity --- Retrieval Ka Formula:}

Do vectors $\vec{a}$ aur $\vec{b}$ ke beech similarity:
\[
\text{similarity}(\vec{a}, \vec{b}) = \frac{\vec{a} \cdot \vec{b}}{\|\vec{a}\| \, \|\vec{b}\|} = \frac{\sum_{i=1}^{n} a_i b_i}{\sqrt{\sum_{i=1}^{n} a_i^2} \cdot \sqrt{\sum_{i=1}^{n} b_i^2}}
\]

\begin{itemize}
\item Result $\in [-1, 1]$. 1 = same direction, 0 = orthogonal, $-1$ = opposite.
\item Sentence embeddings normalize hote hain, to cosine similarity = dot product.
\item ChromaDB isi score par top-k results laati hai.
\end{itemize}
\end{mathbox}

\subsection{all-MiniLM-L6-v2 Model --- Chhota Par Zabardast}

\begin{conceptbox}
\textbf{Model specs:}
\begin{itemize}
\item Name: \texttt{sentence-transformers/all-MiniLM-L6-v2}
\item Architecture: MiniLM (distilled BERT) --- 6 transformer layers.
\item Output dimensions: \textbf{384} (BERT-base 768 hota hai, isne compress kiya).
\item Size: \textasciitilde80MB --- RAM mein load hota hai, GPU ki zaroorat nahi.
\item Speed: CPU par bhi fast (isliye Streamlit app start hote hi load ho jata hai).
\item Quality: 384-dim par bhi semantic search ke liye solid.
\end{itemize}
\end{conceptbox}

\begin{memorybox}
\textbf{Why 384 not 768 or 1536?}

Trade-off: \textbf{quality vs speed/memory}.
\begin{itemize}
\item Higher dims = better nuance, par zyada RAM, slower search.
\item \texttt{all-MiniLM-L6-v2} (384) is project ke scale ke liye perfect: 1000s documents, CPU-only environment, real-time search chahiye.
\item OpenAI \texttt{text-embedding-3-small} (1536 dims) zyada powerful par paid API + network call.
\item Interview answer: ``Mainne quality-speed trade-off socha --- local, free, fast chahiye tha, isliye 384-dim MiniLM.'' 
\end{itemize}
\end{memorybox}

\subsection{Chunking --- PDF Ko Katna Kaise}

\begin{codestorybox}
\begin{lstlisting}[style=pythonstyle, caption={Chunking settings}]
CHUNK_SIZE = 1000        # har chunk ~1000 characters
CHUNK_OVERLAP = 200      # 200 chars overlap pichle chunk se

splitter = RecursiveCharacterTextSplitter(
    chunk_size=CHUNK_SIZE, chunk_overlap=CHUNK_OVERLAP
)
chunks = splitter.split_documents(docs)
\end{lstlisting}

\textbf{RecursiveCharacterTextSplitter ka kaam:}
\begin{itemize}
\item Pehle \textbf{paragraph breaks} par split karta hai (\texttt{\textbackslash n\textbackslash n}).
\item Phir \textbf{sentences} (\texttt{\textbackslash n}, .), phir words, phir characters.
\item Har level par chunk\_size limit check karta hai.
\item Overlap ensure karta hai ki context chunk boundary par na kate.
\end{itemize}
\end{codestorybox}

\begin{warningbox}
\textbf{Chunking ke 3 common galtiyan:}
\begin{enumerate}[label=\arabic*.]
\item \textbf{Too small} (100 chars): har chunk mein adhura context --- retrieval par garbage answers.
\item \textbf{Too big} (5000+): ek chunk mein bahut saari cheezein --- embedding dilute, irrelevant chunks retrieve hote hain.
\item \textbf{Zero overlap}: Sentence beech mein kat jata hai, meaning khota hai.
\end{enumerate}
Is project mein 1000/200 balanced choice hai --- generic documents ke liye.
\end{warningbox}

\subsection{ChromaDB --- Vector Store Ke Andar}

\begin{conceptbox}
\textbf{ChromaDB ki key concepts:}
\begin{itemize}
\item \textbf{Collection:} Vector ka folder (\texttt{research\_docs}, \texttt{past\_research}).
\item \textbf{Embedding function:} Kisi bhi text ko vector mein convert karne wala callable.
\item \textbf{Persistence:} \texttt{persist\_directory} par disk mein save hota hai (SQLite + HNSW index files).
\item \textbf{Similarity search:} Query embedding $\rightarrow$ collection ke saare vectors se cosine compare $\rightarrow$ top-k.
\end{itemize}
\end{conceptbox}

\begin{mathbox}
\textbf{HNSW --- Approximate Nearest Neighbor:}

10000 chunks ho to har query par saare 10000 se compare karna slow hai. HNSW (Hierarchical Navigable Small World) ek \textbf{graph-based index} banata hai:
\begin{itemize}
\item Layers of proximity graphs (top layer = coarse, bottom = fine).
\item Search upper layer se shuru, greedily neeche aata hai.
\item \textbf{Recall} \textasciitilde95\% par \textbf{speed} 10-100x faster.
\end{itemize}
\textbf{Interview line:} ``ChromaDB HNSW index use karti hai --- approximate search se quality ka zyada nuksan nahi, par speed ka bada faida.''
\end{mathbox}

\subsection{Retrieval Quality --- Kya Galt Ho Sakta Hai}

\begin{interviewbox}
\textbf{Q: ``Agar retrieval galat chunks laaye to kya hoga?''}

\textbf{Answer:} Garbage in, garbage out --- RAG system ka output utna hi acha jitna retrieved context. Is project mein 3 safeguards:
\begin{enumerate}[label=\arabic*.]
\item \textbf{Multi-source:} Tavily web results + research\_docs + past\_research teen alag sources se context milta hai.
\item \textbf{Analyst filter:} Analyst sirf \texttt{research\_data[:10]} aur \texttt{rag\_context[:5]} le kar dense brief banata hai --- raw retrieval se pehle hi refine.
\item \textbf{Fact-checker:} Har claim sources se verify hota hai --- hallucinated retrieval bhi pakda jata hai.
\end{enumerate}
\end{interviewbox}

\begin{memorybox}
\textbf{Failure modes table:}

\begin{longtable}{p{4cm} p{5cm} p{4.5cm}}
\toprule
\textbf{Failure} & \textbf{Kya hota} & \textbf{Mitigation} \\
\midrule
Chunk context cut & Adhura answer & Overlap 200 \\
Wrong collection & Galat docs retrieve & Collection names separate \\
Empty collection & Kuch retrieve nahi & try/except $\rightarrow$ 0 stats \\
Embedding mismatch & Garbage similarity & Same model everywhere \\
Query too specific & No good match & Multi-query expansion (researcher) \\
\bottomrule
\end{longtable}
\end{memorybox}
